\documentclass[11pt,a4paper]{article}

\usepackage[utf8]{inputenc}
\usepackage[T1]{fontenc}
\usepackage{geometry}
\usepackage{graphicx}
\usepackage{titlesec}
\usepackage{setspace}
\usepackage{hyperref}
\usepackage{booktabs}
\usepackage{enumitem}
\usepackage{array}
\usepackage{xcolor}
\usepackage{parskip}

\geometry{top=1in, bottom=1in, left=1.25in, right=1.25in}
\hypersetup{colorlinks=true, linkcolor=black, urlcolor=blue, citecolor=black}
\setlength{\parskip}{4pt}
\setlength{\parindent}{0pt}

% Section style
\titleformat{\section}{\large\bfseries}{}{0em}{}[\vspace{-4pt}\rule{\linewidth}{0.5pt}]
\titlespacing{\section}{0pt}{14pt}{6pt}

% Compact list style
\setlist[itemize]{noitemsep, topsep=3pt, leftmargin=*}
\setlist[enumerate]{noitemsep, topsep=3pt, leftmargin=*}

% ============================================================
\begin{document}

% ── TITLE PAGE ──────────────────────────────────────────────
\begin{titlepage}
    \centering

    \includegraphics[width=0.22\textwidth]{iiitl.png}\\[0.8cm]

    {\large \textbf{INDIAN INSTITUTE OF INFORMATION TECHNOLOGY LUCKNOW}}\\[1.2cm]

    {\LARGE \textbf{A Data-Driven Framework for\\[0.25cm]
    Wind Energy Deviation Penalty Optimization}}\\[1.2cm]

    {\normalsize \textit{A project report submitted in partial fulfillment of the requirements}}\\[0.4cm]
    {\normalsize Submitted by}\\[0.35cm]

    {\large \textbf{Tapas Barman}}\\[0.15cm]
    {\normalsize Roll No.: \textbf{MSD24021}}\\[0.15cm]
    {\normalsize MSc.\ Data Science}\\[1.2cm]

    {\normalsize Under the Supervision of}\\[0.35cm]
    {\large \textbf{Dr. Sirsendu Sekhar Barman}}\\[0.15cm]
    {\normalsize Department of Mathematics, IIIT Lucknow}\\[1.8cm]

    {\normalsize March 08, 2026}
\end{titlepage}

% ── TABLE OF CONTENTS (enumerate style) ─────────────────────
\newpage
\thispagestyle{empty}

\begin{center}
    {\large \textbf{Contents}}
\end{center}
\vspace{0.4cm}

\begin{enumerate}
    \item Objective and Introduction
    \item Problem Statement
    \item Proposed Methodology
    \item Dataset and Work Done
    \item Results and Future Work
    \item References
\end{enumerate}

\vfill
\hrule
\vspace{0.3cm}
\begin{center}{\small 1}\end{center}

% ── SIGNATURES PAGE ──────────────────────────────────────────
\newpage
\thispagestyle{empty}

\begin{center}
    {\large \textbf{Signatures}}
\end{center}
\vspace{0.8cm}

\noindent \textbf{Dr. Sirsendu Sekhar Barman}\\
Department of Mathematics\\
Indian Institute of Information Technology, Lucknow\\[1.2cm]
\noindent Signature:\enspace \underline{\hspace{6cm}}

\vspace{1.8cm}

\noindent \textbf{Tapas Barman}\\
Roll No.: MSD24021 \quad MSc.\ Data Science\\
Indian Institute of Information Technology, Lucknow\\[1.2cm]
\noindent Signature:\enspace \underline{\hspace{6cm}}

\vfill
\hrule
\vspace{0.3cm}
\begin{center}{\small 2}\end{center}

% ── CONTENT PAGES ────────────────────────────────────────────
\newpage
\setcounter{page}{3}

% ────────────────────────────────────────────────────────────
\section{Objective and Introduction}

India targets \textbf{500 GW} of renewable energy by 2030, with wind power as a critical pillar. However, wind generation is inherently variable, creating a major grid management challenge.

\textbf{The Scheduling Rule:} All wind farms must submit a \textbf{day-ahead generation schedule} split into \textbf{96 blocks of 15 minutes each}, declared 1.5--2 hours before each slot. Under CERC's \textbf{Deviation Settlement Mechanism (DSM) 2024}, errors beyond 10\% of declared output attract escalating financial penalties.

\textbf{Project Objective:} Build a \textbf{Physics-Informed Predictive Machine Learning (PI-PML)} pipeline that:
\begin{itemize}
    \item Uses physics equations to compute a power baseline from weather forecast data
    \item Applies an XGBoost model trained on real historical data to correct remaining errors
    \item Produces an accurate 96-block day-ahead schedule, minimising CERC DSM penalties
\end{itemize}

% ────────────────────────────────────────────────────────────
\section{Problem Statement}

\textbf{Financial Audit (RSOPL Koppal, Apr 2025 -- Feb 2026):}
\begin{itemize}
    \item Under-production penalties: \textbf{Rs.\ 41.34 Crores}
    \item Average scheduling error: \textbf{5.96 MW} (7.95\% of 75 MW capacity)
    \item Errors are dangerously close to the 10\% DSM penalty threshold
\end{itemize}

\textbf{Root Causes:}
\begin{itemize}
    \item Global NWP weather models (25--30 km grid) cannot resolve micro-climatic wind patterns at a single farm site
    \item Standard tools ignore \textbf{humidity effects} on air density --- during Karnataka's June--September monsoon, wet air reduces power output by 2--4\% compared to dry-air assumptions
    \item Evening hours (6--10 PM) carry a \textbf{2.5$\times$ DSM price multiplier}, making mis-forecasting extremely costly
\end{itemize}

% ────────────────────────────────────────────────────────────
\section{Proposed Methodology}

The PI-PML pipeline operates in two phases:

\textbf{Phase 1 --- Training (Offline):}
\begin{itemize}
    \item SRPC historical generation records are used as ground-truth labels
    \item Physics Engine computes a theoretical baseline $P_\text{physics}$ for each block
    \item Residual = $P_\text{actual}^\text{SRPC}$ $-$ $P_\text{physics}^\text{NWP}$
    \item XGBoost is trained to predict this residual gap
\end{itemize}

\textbf{Phase 2 --- Live Day-Ahead Inference (to be built in End-Semester):}
\begin{itemize}
    \item NOAA GFS weather forecast (wind, temperature, pressure at 100 m) will be fetched every morning
    \item Physics Engine will convert it to a 96-block power profile
    \item Trained XGBoost model will add its learned correction per block
    \item Final corrected schedule will be submitted to CERC for DSM compliance
\end{itemize}

\textbf{Physics Engine --- Three Steps:}
\begin{itemize}
    \item \textbf{Moist Air Density} (Tetens Eq.\ + Modified Gas Law): Corrects for humid air being lighter, fixing monsoon-season over-predictions
    \item \textbf{Wind Height Correction} (Hellman Power Law): Scales 10 m NWP wind to 100 m hub height --- wind is $\approx$32\% faster at 100 m, giving $\approx$2.3$\times$ more power
    \item \textbf{Turbine Power Curve} (Betz Limit, \texttt{windpowerlib}): Maps corrected wind speed and air density to MW output, capped at 59.3\% efficiency
\end{itemize}

% ────────────────────────────────────────────────────────────
\section{Dataset and Work Done}

\textbf{Datasets Identified and Collected:}

\begin{table}[h!]
\centering
\small
\begin{tabular}{@{} p{3.5cm} p{9.7cm} @{}}
    \toprule
    \textbf{Dataset} & \textbf{Details} \\
    \midrule
    SRPC Historical Data &
        Southern Regional Power Committee, Govt.\ of India
        (\url{https://srpc.kar.nic.in}).
        45 weekly ZIP archives, Apr 2025 -- Feb 2026,
        totalling \textbf{29,568 fifteen-minute generation blocks} for
        RSOPL Koppal (75 MW). Used as ground-truth training labels. \\[4pt]
    NOAA GFS Forecast Data &
        National Oceanic and Atmospheric Administration -- Global Forecast
        System (\url{https://nomads.ncep.noaa.gov}).
        Provides wind speed, temperature, and pressure at native
        \textbf{100 m height} (matches turbine hub), updated \textbf{4 times/day},
        free API. Will be the live weather input for day-ahead inference. \\
    \bottomrule
\end{tabular}
\caption{Datasets identified and prepared for the day-ahead pipeline}
\end{table}

\textbf{Work Completed So Far:}
\begin{itemize}
    \item Conducted a full \textbf{10-month CERC DSM financial audit} on SRPC data confirming Rs.\ 41.34 Crores in penalty liabilities
    \item Downloaded, extracted, and merged all 45 SRPC weekly archives into a clean master dataset of 29,568 blocks
    \item Evaluated three NWP weather sources (NASA GEOS FP, NOAA GFS, BFS) and selected NOAA GFS as the primary live input
    \item Built the \textbf{Physics Engine} in Python using \texttt{windpowerlib}: computes moist air density, applies Hellman wind height correction, and maps output through the turbine power curve
    \item Trained an \textbf{XGBoost residual model} on 8 months of SRPC data; validated on the held-out final 2 months
\end{itemize}

% ────────────────────────────────────────────────────────────
\section{Results and Future Work}

\textbf{Training Results (Backtested on 2-Month Holdout, Jan--Feb 2026):}

\begin{table}[h!]
\centering
\small
\begin{tabular}{@{} l c c @{}}
    \toprule
    \textbf{Model} & \textbf{Mean Error (MW)} & \textbf{vs.\ Baseline} \\
    \midrule
    Original declared schedule    & 5.96 & --- \\
    Physics Engine only           & 3.84 & $-$35.6\% \\
    PI-PML (Physics + XGBoost)    & \textbf{1.71} & \textbf{$-$71.2\%} \\
    \bottomrule
\end{tabular}
\caption{Model comparison on held-out 2-month test set}
\end{table}

\begin{itemize}
    \item Error reduced from 5.96 MW to \textbf{1.71 MW} --- a \textbf{71.2\% improvement} over the current schedule
    \item 147 of 192 penalty-triggering blocks corrected to within the 10\% safe zone
    \item Estimated \textbf{Rs.\ 23.5 Lakhs in penalties avoided} in the 2-month test window alone
\end{itemize}

\textbf{Upcoming Deliverables --- Complete Live Day-Ahead Pipeline:}

The trained model and both datasets are in place. The next phase will integrate all developed components into a \textbf{fully automated, production-ready day-ahead forecasting pipeline}. The system will operate as follows:

\begin{enumerate}
    \item \textbf{Automated Weather Ingestion:} Scheduled daily fetch of NOAA GFS 24-hour forecasts for the RSOPL Koppal coordinates, parsing wind speed, temperature, and surface pressure at 100 m using the GRIB filter API
    \item \textbf{Terrain-Aware Wind Correction:} Apply a site-specific roughness correction factor calibrated for Koppal's Deccan Plateau hill terrain, accounting for local surface roughness and orographic flow effects beyond the standard Hellman law
    \item \textbf{Physics Baseline Generation:} Run the full Physics Engine (moist air density, height correction, turbine power curve) to produce a 96-block theoretical power profile for the next day
    \item \textbf{XGBoost Residual Correction:} Apply the trained model to generate a physics-corrected, data-driven generation schedule for each 15-minute block, clipped to valid operational bounds [0, 75 MW]
    \item \textbf{Probabilistic Schedule Output (P10/P50/P90):} Generate low, median, and high confidence bands to allow the plant operator to select a risk-adjusted submission strategy tailored to current weather uncertainty
    \item \textbf{CERC-Formatted Schedule Submission:} Output the final 96-block day-ahead schedule in the standard DSM-compliant format required for CERC submission
    \item \textbf{Financial Impact Validation:} Compute block-by-block DSM penalty and credit comparisons between the PI-PML corrected schedule and the original declared schedule, producing a month-by-month financial savings report
\end{enumerate}

% ────────────────────────────────────────────────────────────
\section{References}

\begin{enumerate}[label={[\arabic*]}, leftmargin=*, itemsep=4pt]
    \item CERC, \textit{Deviation Settlement Mechanism and Related Matters Regulations, 2024}. \url{https://www.cercind.gov.in}
    \item SRPC, \textit{Commercial DSM Accounts --- RSOPL Koppal, 2025--26}. \url{https://srpc.kar.nic.in}
    \item T.\ Chen \& C.\ Guestrin, ``XGBoost: A Scalable Tree Boosting System,'' \textit{ACM SIGKDD}, pp.\ 785--794, 2016.
    \item A.\ Karpatne et al., ``Theory-Guided Data Science,'' \textit{IEEE Trans.\ Knowl.\ Data Eng.}, vol.\ 29, no.\ 10, 2017.
    \item windpowerlib Team, \textit{windpowerlib: Python library for wind turbine modelling}, 2023. \url{https://windpowerlib.readthedocs.io}
    \item NOAA, \textit{Global Forecast System (GFS) Product Description}, 2024. \url{https://www.nco.ncep.noaa.gov}
\end{enumerate}

\end{document}
